\documentclass[11pt]{article}
\usepackage{amsmath,amssymb,amsthm,xcolor,geometry,hyperref,booktabs}
\geometry{margin=1in}
\newcommand{\chg}[1]{\textcolor{blue}{#1}}
\newcommand{\tr}{\operatorname{tr}}
\newcommand{\Res}{\operatorname{Res}}
\newcommand{\arctanh}{\operatorname{arctanh}}
\newcommand{\arccosh}{\operatorname{arccosh}}
\newcommand{\Pf}{\operatorname{Pf}}
\title{Challenge 8: Scalar spectrum in Nieh--Yan gravity \\ \large Problem statement (with revisions marked in blue)}
\date{}
\begin{document}\maketitle
\noindent\textbf{Problem ID:} \texttt{Challenge\_8\_main}\\
\textbf{Primary inferred discipline:} Gravitation, Cosmology \& Astrophysics\\
\textbf{Secondary inferred disciplines:} High Energy \& Nuclear Physics

\section*{Problem setup}
In order to introduce torsion to the system, one can use the first-order formulation of general relativity. We define a local reference frame at each point of the $(3+1)$-dimensional manifold $\mathcal M$, the tetrad $e^A_\mu$, such that the metric can be written as $g_{\mu\nu}=e^A_\mu e^B_\nu\eta_{AB}$, where $\eta_{AB}$ is the flat Minkowski metric on the internal space of coordinates. The internal indices, denoted by the Latin alphabet, also run from 0 to 3 just like the spacetime ones. The metrics $g_{\mu\nu}$ and $\eta_{AB}$ can raise or lower the spacetime and internal tangent-space indices, respectively. The Levi-Civita symbol shall be denoted by $\epsilon_{ABCD}$.

The gravitational action can be reformulated in the first-order form as a function of the tetrad ($e^A$) and spin-connection variables ($\omega^{AB}$). Both of these are 1-forms on the manifold $\mathcal M$. In this formalism, the curvature 2-form is
\begin{equation}
R^{AB}=d\omega^{AB}+\omega^A{}_C\wedge\omega^{CB}.
\end{equation}
Start with the Einstein--Hilbert action ($\mathcal S_{EH}$) in first-order Palatini form, and in first-order Palatini form, add an action term ($\mathcal S_\vartheta$) for a single scalar, $\vartheta$, with an as-yet unspecified potential that depends on $\vartheta$, $V(\vartheta)$. Assume that the scalar $\vartheta$ depends only on time, $\vartheta(t)$, and take $c=1$.

We add a Nieh--Yan action, which we write as
\begin{equation}
S_{NY}=-nf\int d\vartheta\wedge T^A\wedge e_A,
\label{eq:NY}
\end{equation}
where $T^A$ is the torsion two-form
\begin{equation}
T^A=de^A+\omega^A{}_B\wedge e^B .
\end{equation}
We introduce the ansatz for the torsion 2-form:
\begin{align}
T^0&=0,\\
T^i&=h(t)\,e^0\wedge e^i-\phi(t)\,\epsilon^i{}_{jk}\,e^j\wedge e^k .
\label{eq:Ti}
\end{align}
We split the spin connection into ``torsion free'' and ``torsion full'' parts: $\omega^{IJ}=\bar\omega^{IJ}+\tilde\omega^{IJ}$.

Assume a FRW geometry. The scale factor is denoted by $a(t)$, where $t$ is the cosmic time and the Hubble parameter is defined as $H(t)$. In the spatially flat gauge, the metric is given by
\begin{equation}
[g_{\mu\nu}]=a^2(\eta)[\eta_{\mu\nu}+h_{\mu\nu}]=a^2(\eta)\begin{bmatrix}1+2A&-\partial_iB\\-\partial_iB&-\delta_{ij}\end{bmatrix},
\end{equation}
where $\eta$ is conformal time. Hence, the components of the tetrad field $e^A{}_\mu$ are given by
\begin{equation}
e^0{}_0=a[1+A],\quad e^0{}_i=a\partial_i\beta,\quad e^a{}_0=a\delta^{ai}\partial_i\zeta,\quad e^a{}_i=a[\delta_{ia}+\epsilon_{aik}\partial_ks],
\end{equation}
where we have defined $B=\zeta-\beta$ and $s$ is a pseudo-scalar. In the scalar sector, we also have the perturbations
\begin{equation}
h=h(\eta)+\delta h(\eta,\vec x),\quad \phi=\phi(\eta)+\delta\phi(\eta,\vec x),\quad \vartheta=\vartheta(\eta)+\delta\vartheta(\eta,\vec x).
\end{equation}

\section*{Main problem}
\begin{enumerate}
\item $P_{\mathcal R}$ is the curvature power spectrum. Use the values $n=0.5$, $V=\Lambda^4[1-\cos(\vartheta/f)]$, $M_{Pl}=1$, $\Lambda=3.7\times10^{-3}$, $f=1.7$, $a[t=0]=10$, $\vartheta[t=0]=5$ and $\dot\vartheta[t=0]=0$, where $\dot\vartheta=d\vartheta/dt$, $M_{Pl}=\frac{1}{\sqrt{8\pi G}}$ and $G$ is the gravitational constant, to give the value of
\begin{equation}
\frac{P_{\mathcal R}\,(1+3n^2f^2)}{\frac{H^2}{4\pi^2M_{Pl}^2}\big(\frac{H}{\dot\vartheta}\big)^2 2^{2\nu-3}\big|\frac{\Gamma(\nu)}{\Gamma(\frac32)}\big|^2}\times\frac{2AH}{\dot\vartheta\,\delta\vartheta}\times\frac{\beta a\dot\vartheta}{\delta\vartheta}\times\frac{\delta\phi}{nf\,\delta\dot\vartheta-nf\,\dot\vartheta A}
\label{eq:E}
\end{equation}
at horizon crossing at 60 e-folds before inflation ends. $\nu$ is defined as
\begin{equation}
\frac{d^2v}{d\eta^2}+\Big[k^2-\frac{\nu^2-\frac14}{\eta^2}\Big]v=0,
\end{equation}
which is obtained after solving for all the perturbations.
\item What is $\dfrac{\delta\phi}{\delta\dot\vartheta-\dot\vartheta A}$?
\item What is $\dfrac{2AH}{\dot\vartheta\,\delta\vartheta}$?
\end{enumerate}

\section*{\chg{Issues found in the statement and how they are handled}}
\chg{\begin{enumerate}
\item \emph{Unstated normalisations.} The forms of $\mathcal S_{EH}$ and $\mathcal S_\vartheta$ are not written out. The standard first-order forms of the source of this problem [Adshead, Brahma, Das, arXiv:2507.10696, eqs.\ (2.1), (2.3)] are adopted: $\mathcal S_{EH}=-\frac{M_{Pl}^2}{4}\int\epsilon_{ABCD}e^A\wedge e^B\wedge R^{CD}$ and $\mathcal S_\vartheta=\int[\frac12\star d\vartheta\wedge d\vartheta-\frac{1}{3!}\epsilon_{ABCD}e^Ae^Be^Ce^D\,V]$, with signature $(+,-,-,-)$, $\epsilon_{0123}=1$, and $\epsilon^i{}_{jk}$ in (\ref{eq:Ti}) the three-dimensional symbol $\epsilon_{ijk}=\epsilon_{0ijk}$. Only the \emph{relative} signs of $\mathcal S_{EH}$, $\mathcal S_\vartheta$ and $S_{NY}$ matter physically: the sign of the Einstein--Hilbert term relative to the scalar kinetic term is fixed by requiring positive energies, and then the sign of $\phi/(nf\dot\vartheta)$ is fixed by (\ref{eq:NY}); with the conventions above $\phi=+nf\dot\vartheta/M_{Pl}^2$. Item 1 of the main problem is independent of this sign.
\item \emph{The factor $(1+3n^2f^2)$ in (\ref{eq:E}).} The kinetic coefficient of $\vartheta$ produced by integrating out torsion is $1+6n^2f^2/M_{Pl}^2$, not $1+3n^2f^2$; the factor is nevertheless kept exactly as printed, and the answer is computed for the expression as written. (If $1+6n^2f^2$ had been intended, the value of (\ref{eq:E}) would be $5.335$ instead of $3.1675$; see the solution.)
\item \emph{Numbering.} The three items were all labelled ``1.'' in the original; they are relabelled 1--3. The tetrad component $e^a{}_i=a[\delta_{ia}+\epsilon_{aik}\partial_ks]$ is read as $e^a{}_i=a[\delta^a_i+\epsilon_{aik}\partial_ks]$.
\item \emph{Role of the numerical data.} The three requested quantities are exact constants (independent of $H$, $\dot\vartheta$, $\nu$, $a$, $\Lambda$ and the mode amplitudes), so $\Lambda$, $a[0]$, $\vartheta[0]$, $\dot\vartheta[0]$ only serve to guarantee that the evaluation point (horizon crossing 60 e-folds before the end of inflation, in the rolling phase with $\dot\vartheta\neq0$) exists; this is verified in the solution.
\end{enumerate}}
\end{document}